\documentclass[11pt,a4paper]{article}
\usepackage[margin=23mm]{geometry}
\usepackage[T1]{fontenc}
\usepackage[utf8]{inputenc}
\usepackage{lmodern,amsmath,amssymb,amsthm,longtable,array,booktabs,xurl,hyperref}
\hypersetup{colorlinks=true,linkcolor=blue,urlcolor=blue}
\setlength{\parindent}{0pt}
\setlength{\parskip}{6pt}
\setlength{\emergencystretch}{3em}
\newcommand{\code}[1]{\nolinkurl{#1}}
\newcommand{\Tr}{\operatorname{Tr}}
\newcommand{\E}{\mathbb E}
\newcommand{\Law}{\operatorname{Law}}
\newcommand{\PhiB}{\operatorname{Phi}}
\newtheorem{proposition}{Propozycja}
\title{FT1536 / S05: kontynuacja rozpoznania QROM\\
\large Literatura, mapa eksportów i warunkowa kompozycja pełnych odpowiedzi}
\author{Notatka robocza dla Niirmata}
\date{10 października 2026 / CONTINUATION\_002}
\begin{document}
\maketitle

\textbf{Status: rozpoznanie literatury oraz dowód tekstowy warunkowego lematu.}
Bez replayu Lean/Sage, kernelowej formalizacji i niezależnego odbioru.
Bez twierdzenia o bezpieczeństwie rzeczywistego FT1536. Rozwijamy S05 z
ROADMAP; nie zmieniamy znaczenia istniejących statusów T12.1 ani B20.

\section{Co zmienia ten krok}
Punktem wyjścia jest zachowany pakiet \code{FT1536-S05-QROM-RECON-001}.
Uzupełniono brakujący odczyt transformacji Unruha i sprawdzono dokładne
twierdzenia KLS, DFMS/DFM oraz adaptacyjnego reprogramowania.
Wynik pozytywny jest węższy od pełnego QROM: skonstruowano wspólny kanał
dla dwóch \emph{przeprogramowanych} gier oraz wyprowadzono adaptacyjny
rachunek drugiego momentu, gdy odpowiedzi pochodzą z niezależnego,
klasycznego samplera o określonym niżej interfejsie.

To przesuwa pierwszy konkretny obowiązek do \textbf{Q-JOINT-INSTANCE}:
trzeba dostarczyć efektywny sampler pełnej odpowiedzi i jego lokalny
certyfikat. Co więcej, kwantowa legalność argumentu stanu ma prostą
warunkową realizację: istniejący \code{LocalJointCertificate} jest
jednolity po \emph{wszystkich} stanach, więc do \code{S.code} można zawsze
podać publiczny \code{initial}. To nowy sampler, nie utożsamienie
kwantowej historii z tablicą klasycznego ROM.
Q-TARGET, rzeczywiste prawo Sign i zasoby nadal są otwarte.

\textbf{Ważne rozróżnienie.} Klasyczne zapytania Sign i kwantowe zapytania
hasza są standardowym modelem QROM dla podpisów, a nie same w sobie
osłabieniem tego modelu [KLS, U15]. Mniejszym wynikiem jest tutaj warunkowy
lemat symulacji. Dostęp w superpozycji do Sign byłby osobnym celem.

\textbf{Baza źródłowa:}
\code{2d2cdf52304cee9cf513487d4296392262df22f4}.
Lokalny HEAD przy odczycie:
\code{eb9836cc9138e89d1f0cbcd052a8a58e32caf94f}.
Wszystkie 24 identyfikatory blobów poprzedniego manifestu sprawdzono lokalnie;
23 podane SHA-256 również się zgadzają. Pliki formalne i C użyte poniżej
są takie same w bazie i w odczytanym drzewie roboczym. Zmiany czterech
dokumentów odnotowano w \code{INPUT_RECEIPT.json}; nie traktujemy ich jako
aktualizacji historycznego snapshotu. Pięć hashy oryginalnego pakietu jest zgodnych.

\section{Dokładne liście literaturowe}
Sprawdzono treść pod wskazanymi adresami 10 października 2026.
Numery stron poniżej są numerami drukowanymi prac, nie indeksem PDF.
Nie wykonano niezależnego audytu całych publikacji. PDF-ów nie pobierano
do repozytorium; adres wersji i data są pinem bibliograficznym, nie hashem bajtów.
Nie formułujemy roszczenia nowości ani kompletności przeglądu literatury.

\subsection*{KLS: konkretny Fiat--Shamir}
Kiltz--Lyubashevsky--Schaffner, \emph{A Concrete Treatment of Fiat-Shamir
Signatures in the Quantum Random-Oracle Model}, zaakceptowany manuskrypt
20 II 2018, DOI 10.1007/978-3-319-78372-7\_18 [KLS].
Definicja 2.5, rysunek 8 i twierdzenia 3.3--3.4 (s.~9, 11, 14).
Twierdzenie 3.3 zakłada $\varepsilon_{zk}$-perfect naHVZK i $\alpha$ bitów
min-entropii zobowiązania; daje
\[
 \operatorname{Adv}_{CMA}(A)\le\operatorname{Adv}_{NMA}(B)
       +q_s2^{1-\alpha}+\kappa_m q_s\varepsilon_{zk},
 \qquad T_B=T_A+\kappa_m q_Hq_s.
\]
Twierdzenie 3.4 wymaga dodatkowo trybu lossy i jego statistical soundness;
zastępuje NMA przez LOSS z kosztem $8(q_H+1)^2\varepsilon_{ls}$.
CUR dotyczy mocnej niepodrabialności.
NaHVZK wymaga m.in. jednostajnego wyzwania w nieabortującej symulacji.
Rysunek 8 odświeża zobowiązanie przy każdej zewnętrznej próbie.
\textbf{Zastosowanie: warunkowe.} Słownik kandydata to $W=r$, $c=H(r\Vert m)$,
$Z=s$. Należałoby włożyć całą pętlę 16 prób do odpowiedzi jednego protokołu
i osobno obsłużyć widoczne porażki. Ani wymagane naHVZK, ani tryb lossy
dla rzeczywistego rozkładu kluczy nie wynikają z lokalnego momentu FT1536.

\subsection*{DFMS19 i DFM20: dwa różne kontrakty}
Don--Fehr--Majenz--Schaffner, \emph{Security of the Fiat-Shamir Transformation
in the Quantum Random-Oracle Model}, arXiv:1902.07556v4, 27 VII 2020 [DFMS19].
Twierdzenie 22 (s.~14) jest asymptotycznym wynikiem sEUF-CMA: wymaga trudnej
relacji, superwielomianowej przestrzeni wyzwań, quantum proof of knowledge
według definicji 14 oraz naHVZK, entropii i CUR z KLS.
\textbf{Zastosowanie: brak instancji.} \code{accepted_extracts} wydaje krótki
przeciwobraz pojedynczego wyzwania; nie dowodzi wymaganej wiedzy świadka
relacji identyfikacyjnej. Sam rozmiar $R_q$ dla jednego profilu nie zastępuje
rodziny asymptotycznej. Ten wynik nie jest uzasadnieniem automatycznego
podstawienia nowego przeciwnika do \code{concrete_euf_cma_to_mt_isis}.

Don--Fehr--Majenz, \emph{The Measure-and-Reprogram Technique 2.0:
Multi-Round Fiat-Shamir and More}, arXiv:2003.05207v3, 7 III 2022 [DFM20].
Twierdzenie 2 (s.~5): dla skończonych niepustych $X,Y$, losowego
$H:X\to Y$ i algorytmu o $q$ zapytaniach, dwustopniowy symulator wybiera
$x$ przed otrzymaniem jednostajnego $\Theta\in Y$. Dla każdego $x_0$
i dopuszczalnego predykatu $V$:
\[
 \Pr[x=x_0\land V(x,\Theta,z):S^A]
 \ge\frac{\Pr[x=x_0\land V(x,H(x),z):A^H]}{(2q+1)^2}.
\]
Uwzględniony jest końcowy wynik jako miejsce pomiaru; nie ma dodatkowego
błędu addytywnego z wcześniejszej wersji techniki. Konwencja dopuszcza
kontrolowane zapytania. \textbf{Zastosowanie: kandydat do Q-TARGET}, z
$X$ równym dziedzinie ramek, $Y=R_q$, a $V$ opartym na Verify.
Trzeba wliczyć zapytania symulatora i rozwiązać jego interakcję z Sign.
Twierdzenie nie dostarcza publicznego samplera podpisu.

\subsection*{U15: transformacja, nie certyfikat starego formatu}
Dominique Unruh, \emph{Non-Interactive Zero-Knowledge Proofs in the Quantum
Random Oracle Model}, EUROCRYPT 2015, DOI 10.1007/978-3-662-46803-6\_25;
sprawdzona wersja konferencyjna IACR [U15]. Definicja 8, rysunek 1,
twierdzenia 10, 13, 18 i wniosek 14 (s.~6, 12--13, 20, 22, 24).
Konstrukcja używa wielu transkryptów i dwóch wyroczni $G,H$.
Twierdzenie 10 wymaga HVZK, superlogarytmicznej min-entropii odpowiedzi
warunkowo względem jej stanu/wejścia oraz odpowiedniej entropii kolizyjnej
krotki haszowanej. Twierdzenie 13 wymaga special soundness oraz
superlogarytmicznych $t\log m$, entropii symulowanej krotki i długości odpowiedzi.
Wniosek 14 daje istnienie odpowiedniego NIZK po dopuszczonych modyfikacjach
protokołu. Twierdzenie 18 buduje silnie niepodrabialny podpis z ZK,
simulation-sound online-extractability i generatora trudnych instancji
z definicji 16. \textbf{Zastosowanie do niezmienionego FT1536: nie.}
To inny format i inny algorytm weryfikacji, nie brakująca przesłanka
dla obecnego Sign. Wcześniejszy O2H Unruha nie jest tym twierdzeniem.

\subsection*{AR: opłata za zmianę istniejącej wyroczni}
Grilo--Hövelmanns--Hülsing--Majenz, \emph{Tight adaptive reprogramming
in the QROM}, arXiv:2010.15103v2, 30 X 2020 [AR], twierdzenie 1,
rysunek 2, s.~6--7. Wybrane adaptacyjnie rozkłady losują miejsca
przeprogramowania; nowa wartość jest świeża i jednostajna. Koszt rozróżnienia
jest ograniczony przez sumę
$\sqrt{\widehat q_i p^{(i)}_{max}}+\widehat q_i p^{(i)}_{max}/2$,
gdzie $\widehat q_i$ liczy wcześniejsze zapytania, a $p^{(i)}_{max}$ jest
oczekiwaną największą masą marginalnego rozkładu miejsca.
\textbf{Zastosowanie:} dla nowego idealnego $r\in\{0,1\}^{320}$ i już
ustalonego $m$, $p_{max}=2^{-320}$. Wliczamy także klasyczne odczyty hasza
przez Sign i końcowego weryfikatora. Dla $Q=q_H+q_s+1$ wystarczy
\[
 \delta_{AR}=\min\{1,\ q_s(\sqrt{Q\,2^{-320}}+Q\,2^{-320}/2)\}.
\]
To podstawienie wymaga idealnego nonce, jednego odczytu hasza na Sign
i braku innych dostępów; rzeczywisty PRG nie spełnia go na mocy samej długości.

\subsection*{BM26: dokładna granica gotowej drogi Falcona}
Beckmann--Majenz, \emph{Quantum Oracle Distribution Switching and
Applications to Falcon and Ring Signatures}, arXiv:2602.16268v3,
9 IX 2026 [BM26]. Twierdzenie 5, rysunek 9, dodatek E.1:
gotowy wynik dotyczy CoreFalcon$^+$, który odświeża nonce i wyzwanie w pętli.
Rysunek odróżnia go od wariantu ze stałym nonce. W FT1536 pozostaje jeszcze
cap, inne prawo samplera, norma, rozkład kluczy i emisja bajtów.
Twierdzenie 4 (s.~14) daje przeszkodę dla uniwersalnego przenoszenia
powodzenia między całymi wyroczniami o różnych rozkładach wyjść,
przy stracie niezależnej od rozmiaru dziedziny, nawet dla dwóch zapytań.
\textbf{Wniosek:} nie wolno podmienić wszystkich wartości kwantowej
wyroczni wyłącznie na podstawie lokalnej dywergencji. Nie wyklucza to
poniższej zamiany skończonej liczby klasycznych losowań po opłaconym
reprogramowaniu. Nie przenosimy twierdzenia 5 przez podstawienie parametrów.

\section{Mapa naszych lematów}
Skrót R oznacza \code{proofs/ft1536/development/T12_1/run2/formal/}.
Skrót F oznacza \code{proofs/ft1536/stages/FT1536_MATH_EUFCMA_MTISIS_RUN_001/formal/FT1536/}.
Oceny poniżej dotyczą odczytanych eksportów w przypiętej bazie;
nie są audytem wszystkich dowodów ani nadaniem statusu REVIEWED.

\small
\begin{longtable}{@{}p{.34\linewidth}p{.61\linewidth}@{}}
\toprule Eksport / konstrukcja & Co przeżywa i co trzeba dowieść \\\midrule
\endhead
F/Relation: \code{extraction_equation}, \code{accepted_extracts}
& \textbf{Zachować bez zmiany tezy.} Deterministyczny Verify implikuje
ShortPreimage dla zmierzonego klasycznego wyniku. Rozkład i osadzenie celu
są osobnymi obowiązkami.\\
R/SignLayerSupport: \code{signBodyOf_mass_none}, \code{cap_le_of_pointwise}
& \textbf{Zachować warunkowo.} Klasyczne masy prób i końcowej emisji nadal
mają sens. Warunki zgodności z realnym prawem C i wszystkimi porażkami
nie znikają.\\
R/AttemptPointwise: \code{AttemptShape}, \code{Sandwich}
& \textbf{Zachować jako przesłanki.} Model przeciwnika nie zmienia
nierówności punktowych; nie dostarcza ich instancji źródłowej.\\
R/Run2/LawBinding: \code{honest_joint_body}
& \textbf{Zachować definicję} $P_h(c,o)=|R_q|^{-1}\,\mathrm{signBody}_h(c)(o)$,
z pełnym $o$, także \code{none}.\\
R/Run2/LawBinding: \code{LocalJointCertificate}
& \textbf{Zachować certyfikat, zmienić sposób użycia.} Kierunek $J\ll P$,
$\sum J^2/P\le1+e$ jest właściwy. Uniwersalny kwantyfikator po stanie
pozwala użyć \code{S.code} zawsze z \code{initial}; patrz specjalizacja poniżej.\\
F/EventTransfer: \code{event_quadratic}, \code{phi_bound}
& \textbf{Zachować rachunek.} Dla testu $v\in[0,1]$ i wspólnych stanów
warunkowych działa dowód z RECON-001; nie dla samych marginałów odpowiedzi.\\
R/Run2/Games: \code{parse_frame}, \code{framing_injective}, \code{nonce_card}
& \textbf{Zachować.} To algebra bajtów i liczność. Liczność $2^{320}$ nie
jest twierdzeniem o entropii rzeczywistego generatora.\\
R/Run2/Games: \code{Program}, \code{ClassicalAdversary}, \code{recordedHash}
& \textbf{Zastąpić semantykę.} Klasycznej gałęzi po każdej odpowiedzi hasza
nie wolno utożsamić z zapytaniem w superpozycji. Nie mierzymy rejestrów
kwantowych tylko po to, aby sporządzić log.\\
R/Run2/StoppedComparison: \code{programComparison},
\code{stopped_game_event_bound}
& \textbf{Stare dowody nie obejmują QROM.} Poniższa propozycja daje osobny
dowód tekstowy analogicznej potęgi dla zawężonego interfejsu; nie jest
kwantową instancją istniejącego typu Lean.\\
R/Run2/Games: \code{hashTargets}, \code{finishSim};
R/Run2/ConcreteReduction: \code{concrete_euf_cma_to_mt_isis}
& \textbf{Nie przenosić.} Lista $q_H+1$ kolejnych celów zakłada klasyczne
zapytania. Q-TARGET wymaga nowej konstrukcji, straty i kosztu.\\
R/Run2/ConcreteReduction: \code{exact_collision_parameter}
& \textbf{Zachować tylko klasyczne znaczenie.} Nie jest kwantowym boundem
niezauważalności reprogramowania. Kandydat QROM to osobne $\delta_{AR}$.\\
Realny seed/PRG, S04 i Resources
& \textbf{Otwarte.} Potrzebne bezpieczeństwo PRG przeciw przeciwnikowi
kwantowemu, publiczny bit-output SHAKE/H2P i kwantowy koszt reduktora.\\
\bottomrule
\end{longtable}
\normalsize

\textbf{„Nie przenosić” nie znaczy „fałszywe”.} Klasyczne twierdzenia pozostają
twierdzeniami o swoich typach. Brak kwantowej semantyki jest luką w nowym
dowodzie, nie kontrprzykładem FT1536. Natomiast utożsamienie samych marginałów
z całym stanem jest fałszywą ogólną regułą, jak pokazuje przykład z RECON-001.

\section{Mały wynik: wspólny kanał i adaptacyjna kompozycja}
\textbf{Dowód autorski tekstowy, samosprawdzenie; bez roszczenia nowości.}
Ustalmy skończoną dziedzinę wejść hasza (jawny limit długości), kodowanie
$R_q$, klasyczny klucz publiczny $h$ i skończony horyzont obliczeń.
Wyrocznia działa koherentnie jako
$U_H\lvert x,y\rangle=\lvert x,y\mathbin{\oplus}\mathrm{enc}(H(x))\rangle$;
wartość $H(x)$ jest jednostajna w $R_q$, nie we wszystkich bitstringach
rejestru kodowania. Przeciwnik jest kwantowy, komunikacja Sign jest klasyczna,
a wynik końcowy jest klasyczny. Nie dodajemy dostępu do Sign w superpozycji.

\subsection*{Dopuszczalny interfejs}
Po historii $t$ i klasycznym żądaniu $m$ serwer losuje niezależny,
jednostajny 40-bajtowy nonce $r$. Następnie losuje całą parę $x=(c,o)$
według $P_{t,m,r}$ albo $J_{t,m,r}$ na tym samym skończonym zbiorze.
Obie gry mają tę samą listę złożonych wiadomości, także przy porażce.
W zastosowaniu modelowym
\[
 P_{t,m,r}(c,o)=P_h(c,o)
   =\frac{1}{|R_q|}\,\mathrm{signBody}_h(c)(o),
 \quad o\in\operatorname{Option}(\operatorname{BoxVec}).
\]
Rozkład obejmuje cap16 i końcową porażkę emisji w modelowym prawie.
PRE\_ABORT, UB, timing i inne wyjścia realnego API nie zostały przez tę
definicję utożsamione z \code{none}; ich realizacja wymaga osobnego bridge'u.

Sampler używa świeżej prywatnej klasycznej losowości. Jego rozkład zależy
wyłącznie od $h,m,r$ i dozwolonej klasycznej historii $t$, a \emph{nie}
od ukrytej tablicy $H$, nieobserwowanych zapytań, prywatnego rejestru
przeciwnika ani od uzyskanego przez dodatkowe zapytanie starego $H(r\Vert m)$.
Jeżeli sampler zachowuje prywatny stan, jego zmiana musi być częścią $x$
i certyfikatu albo deterministyczną funkcją zarejestrowanej historii.
Sam certyfikat marginalnej pary $(c,o)$ tego nie pokrywa.

Po losowaniu obie gry wykonują \emph{tę samą} operację:
nadpisują punkt $r\Vert m$ wartością $c$, zwracają $(r,o)$ i kontynuują.
Poprzednie nadpisania są obsługiwane identycznie (ostatni zapis wygrywa).
Brak testu „czy przeciwnik już zapytał ten punkt”.
Każda dodatkowa obserwowalna gałąź ma być wspólna lub uwzględniona w $x$.

\begin{proposition}[Kompozycja dla programowanych pełnych odpowiedzi]
Załóżmy opisany interfejs, najwyżej $q_s$ wywołań Sign oraz, dla każdej
dopuszczalnej historii, wiadomości i nonce,
\[
 J_{t,m,r}\ll P_{t,m,r},\qquad
 \sum_{x:P_{t,m,r}(x)>0}\frac{J_{t,m,r}(x)^2}{P_{t,m,r}(x)}\le1+e,
 \qquad e\ge0.
\]
Dla dowolnego końcowego zdarzenia, niech $a$ będzie jego prawdopodobieństwem
w grze programowanej z $P$, a $b$ w grze programowanej z $J$. Wtedy
\[
 (a-b)^2\le D\,a(1-a),\qquad a\le\PhiB(D,b),\qquad
 D=(1+e)^{q_s}-1,
\]
gdzie $\PhiB(D,b)=(2b+D+\sqrt{D^2+4Db(1-b)})/(2(1+D))$.
\end{proposition}

\begin{proof}
Najpierw jedna odpowiedź. Dla ustalonej historii obie gry mają ten sam
znormalizowany stan $\rho$ rejestrów przeciwnika i mechanizmu wyroczni.
Nie zakładamy, że $\rho$ jest niezależny od $H$.
Świeże losowanie $x$ jest jednak niezależne od tych rejestrów warunkowo
względem historii. Dla wspólnego kanału nadpisania i wydania odpowiedzi
$\mathcal E_{t,m,r,x}$ stan po kroku wynosi odpowiednio
\[
 \sum_xP_{t,m,r}(x)\,\mathcal E_{t,m,r,x}(\rho),\qquad
 \sum_xJ_{t,m,r}(x)\,\mathcal E_{t,m,r,x}(\rho).
\]
Są to rzeczywiście wspólne stany warunkowe. Kanał może być opisany przez
ukrytą klasyczną funkcję $H$ i rejestry przeciwnika; nadpisanie jest kanałem
z zachowaniem śladu, nie pomiarem zapytań przeciwnika. Implementacja
interfejsu do nadpisanej wyroczni używa koherentnego testu równości wejścia
z zapamiętanym punktem. Nie odczytuje całej funkcji $H$.

Rozwijamy teraz analizę w skończone drzewo \emph{istniejących} klasycznych
wyników: żądań Sign, odpowiedzi, nonce, prywatnych par $(c,o)$ i wszystkich
rzeczywiście wykonanych pomiarów klasycznych. Rejestrujemy prywatne
klasyczne wartości tylko na potrzeby analizy, bez ich ujawniania.
Nie dodajemy pomiarów zapytań w superpozycji.
Rozszerzona historia analityczna nie poszerza dostępu samplera: jego
parametrem nadal jest wyłącznie dozwolona projekcja tej historii.
W korzeniu stan warunkowy jest wspólny. W kroku losowania różni się tylko
skalarna masa $P(x)$ lub $J(x)$, więc po ustaleniu $x$ stan znormalizowany
pozostaje wspólny. W kroku wspólnego kwantowego instrumentu prawdopodobieństwo
każdego klasycznego wyniku i kolejny stan warunkowy są wobec tego wspólne.
To dowodzi indukcyjnie istnienia wspólnych stanów przy całej historii,
także gdy następna wiadomość zależy od wcześniejszych wyników kwantowych.

Niech $p(t),j(t)$ będą masami historii; na gałęziach o dodatniej masie $p$
piszemy $L(t)=j(t)/p(t)$. Zachodzi $j\ll p$ przez indukcję.
Wspólne rozgałęzienie o prawdopodobieństwach $k(y\mid t)$ zachowuje
$\sum_t p(t)L(t)^2$, bo $\sum_y k(y\mid t)=1$.
Krok różniącego się losowania spełnia dokładnie
\[
 \sum_{t,x}\frac{(j(t)J_t(x))^2}{p(t)P_t(x)}
 =\sum_t\frac{j(t)^2}{p(t)}
       \sum_{x:P_t(x)>0}\frac{J_t(x)^2}{P_t(x)}
 \le(1+e)\sum_t\frac{j(t)^2}{p(t)}.
\]
Zera masy pomijamy, co jest legalne dzięki absolutnej ciągłości.
Nonce jest wspólnym rozgałęzieniem. Wcześniejsze zakończenia dopełniamy
fikcyjnymi jednoelementowymi losowaniami; mają czynnik 1.
Stąd dla końcowego drzewa $\E_pL=1$ i $\E_pL^2\le(1+e)^{q_s}$.

Dla wspólnego końcowego stanu $\tau_t$ i efektu pomiarowego $0\le M\le I$
niech $v(t)=\Tr(M\tau_t)\in[0,1]$. Mamy $a=\E_pv$, $b=\E_pLv$,
a z Cauchy'ego--Schwarza i $v^2\le v$:
\[
 (b-a)^2
 =\bigl(\E_p[(L-1)(v-a)]\bigr)^2
 \le (\E_pL^2-1)\E_p(v-a)^2
 \le D\,a(1-a).
\]
Górny pierwiastek $(1+D)a^2-(2b+D)a+b^2\le0$ daje podane $\PhiB$.
\end{proof}

Nie utożsamiamy tego dowodu z jego instancją dla istniejącego
\code{Sampler.code}. W szczególności tabela odczytanych klasycznych
hashy z \code{Games.State} nie jest legalnym zapisem wszystkich zapytań QROM.

\subsection*{Specjalizacja istniejącego certyfikatu do stanu początkowego}
Istnieje ważne ułatwienie względem pierwotnego rozpoznania.
Zapisany w \code{Run2/LawBinding.lean} certyfikat kwantyfikuje po
\emph{każdym} $h,st,m,r$, bez warunku osiągalności $st$.
Dla dowolnego $S$ i takiego certyfikatu zdefiniujmy nowy sampler
\[
 S^0(h,m,r;u):=S.\mathrm{code}(h,\mathrm{initial},m,r;u),
 \qquad u\leftarrow\mathrm{Uniform}(\mathrm{Fin}(S.\mathrm{bits})\to\mathrm{Bool}).
\]
Wtedy $J^0_{h,m,r}=\mathrm{samplerLaw}(S,h,\mathrm{initial},m,r)$ i bez
nowego oszacowania, przez podstawienie \code{st := initial}, otrzymujemy
\[
 J^0_{h,m,r}\ll\mathrm{freshHonest}(h),\qquad
 \sum_x\frac{J^0_{h,m,r}(x)^2}{\mathrm{freshHonest}(h)(x)}\le1+e.
\]
W sumie stosujemy tę samą konwencję pomijania zer mianownika.
Argument \code{initial} jest publiczną stałą; \code{S.code} jest funkcją
podanych klasycznych argumentów i świeżej taśmy. Nie ma dostępu do ukrytego
$H$ ani pamięci przeciwnika. $S^0$ spełnia więc warunek niezależności
powyższej propozycji i ma tę samą liczbę losowych bitów.

To \textbf{wyprowadzenie z istniejącej warunkowej przesłanki}, a nie
konstrukcja świadka $S$ ani dowód istnienia certyfikatu.
Nie twierdzimy, że rozkład $S^0$ jest równy rozkładowi dawnego
\code{signSim S} na rzeczywistym stanie: równość nie jest potrzebna,
bo definiujemy nową grę $Q_J$. Efektywna implementacja $S^0$ i jej
konkretny koszt wymagają kodu dla $S$; sam typ funkcji Lean tego nie daje.
Jeśli przyszły certyfikat zostanie osłabiony do stanów osiągalnych,
trzeba ponownie sprawdzić legalność podstawienia \code{initial}.

\subsection*{Połączenie z nieprzeprogramowanym modelem idealnym}
Niech $Q_0$ losuje nonce, odczytuje istniejące $H(r\Vert m)$ i wydaje
odpowiedź według modelowego \code{signBody}. Niech $Q_P$ zamiast tego
losuje $c$ jednostajnie, nadpisuje ten punkt i wydaje odpowiedź tego samego
prawa warunkowego. Reprogramowanie poprzedza obliczenie odpowiedzi.
W grze porównawczej AR odbiorca uzyskuje nonce, wykonuje jeden klasyczny
odczyt aktualnej wyroczni i dopiero losuje odpowiedź. Może znać sekret
na potrzeby porównania dwóch uczciwych gier; nie jest to publiczny reduktor.
Nie pyta o nową wartość przed wykonaniem reprogramowania.

Dla tego modelu podstawienie do [AR] z poprzedniej sekcji daje
$|\Pr[Q_0]-\Pr[Q_P]|\le\delta_{AR}$.
Z propozycji, po wprowadzeniu $Q_J$ z legalnym publicznym samplerem,
otrzymujemy \textbf{warunkowy lemat symulacji}
\[
 \boxed{\Pr[Q_0]\le
 \min\{1,\delta_{AR}+\PhiB((1+e)^{q_s}-1,\Pr[Q_J])\}.}
\]
Licznik $Q=q_H+q_s+1$ obejmuje końcowy odczyt weryfikatora.
Gdy dany model ma więcej odczytów, trzeba zwiększyć $Q$.
Powtórzenia nonce są dopuszczone przez grę reprogramowania; nie używamy
nieuzasadnionego testu świeżości względem superpozycyjnych zapytań.
Liczba reprogramowań wynosi tu najwyżej $q_s$, a nie $16q_s$:
całe ciało jednego Sign, także jego porażka, jest jednym losowaniem.

\textbf{Granica:} $Q_0$ jest modelem współczynnikowym z idealnym nonce
i direct-output hashem. Nie wykazano jego równości z rzeczywistym Sign.
Nie ograniczyliśmy jeszcze $\Pr[Q_J]$ przez trudność ISIS.
To nie jest dowód EUF-CMA ani deklaracja liczby bitów bezpieczeństwa.

\section{Samosprawdzenie i następny pakiet}
Wykonano samosprawdzenie argumentu, nie niezależny odbiór.
Werdykt dla zastosowania do FT1536: \textbf{conditional on a named input}.
Graf zależności: legalność losowania i kanału $\to$ wspólne stany
$\to$ iloczyn momentów $\to$ nierówność Phi; osobno [AR] $\to$
przejście od $Q_0$ do $Q_P$. Żadna z tych strzałek nie dowodzi Q-JOINT.

\begin{itemize}
\item Sprawdzono kierunek $J/P$, normalizację, zera masy, przypadki
$q_s=0$ i $e=0$, zakończenie po porażce oraz adaptacyjny wybór wiadomości.
\item Warunkowa wspólność dotyczy $Q_P,Q_J$. Nie założono jej bez dowodu
dla $Q_0,Q_P$; tam płacimy osobny błąd reprogramowania.
\item Wykryta granica: równość marginałów $(c,o)$ przy różnych ukrytych
korelacjach nie wystarcza. Stan prywatny samplera musi zostać objęty kontraktem.
\item Bibliografia: źródła pierwotne, wskazane wersje i fragmenty.
Brak twierdzenia o nowości; brak skanowania całej literatury.
\item Brak nowych obliczeń Sage i kompilacji Lean: dowód ma status tekstowy.
Kompilacja dokumentu potwierdza skład, nie poprawność matematyczną.
\end{itemize}

\textbf{Następny krok S05-Q-JOINT-INSTANCE.}
Zidentyfikować konkretny efektywny $S$ i jego \code{LocalJointCertificate},
a następnie zastosować $S^0$ ze stałym \code{initial}. Podać dokładną
dystrybucję pełnej porażki i koszt. Jeżeli istnieje tylko warunkowy typ
certyfikatu, skonstruować jego świadka albo wskazać brakujący eksport.
Nie wymagać już rekonstrukcji historii kwantowych zapytań: specjalizacja
do \code{initial} usuwa tę potrzebę na poziomie obecnej przesłanki.

Potem można formalizować wspólny kanał i drzewo kompozycji oraz osobno
konstruować Q-TARGET przez measure-and-reprogram. Jeszcze mniejszym,
odrębnym celem Q-TARGET jest przypadek $q_s=0$ (bez wyroczni podpisu),
z pojedynczym wyzwaniem i \code{accepted_extracts}. Nie odhaczamy go
na podstawie samego wpisania współczynnika $(2q+1)^2$.

\section{Odnośniki do sprawdzonych wersji}
\begin{description}
\item[KLS] \url{https://pure.uva.nl/ws/files/45974858/916.pdf}
\item[DFMS19] \url{https://arxiv.org/pdf/1902.07556v4}
\item[DFM20] \url{https://arxiv.org/pdf/2003.05207v3}
\item[U15] \url{https://www.iacr.org/archive/eurocrypt2015/90560105/90560105.pdf}
\item[AR] \url{https://arxiv.org/pdf/2010.15103v2}
\item[BM26] \url{https://arxiv.org/pdf/2602.16268v3}
\end{description}
Pakiet roboczy: \code{continuation_002/}. Piny wejść i stan dowodowy
zapisano osobno w \code{INPUT_RECEIPT.json} oraz \code{CHECKPOINT.json}.
Oryginalny RECON-001 pozostaje niezmieniony. Brak importu do stages,
commita i push; bieżący materiał pozostaje w przydzielonym W do odbioru.
\end{document}
